\section{Results and Discussion}\label{results-and-discussion}

Chapter 7 established which model performed best. This chapter shifts to interpretation. The point is not to repeat the evaluation table, but to explain what the retained findings mean for residential valuation in Metro Cebu. Read as a whole, the results point to a market that is polycentric, strongly segmented by property type, and only partly captured by simple counts of nearby amenities.

\subsection{Overview of Findings}\label{overview-of-findings}

Three points stand out. First, the strongest locational signals were not limited to the old metropolitan core. Property-type dummies occupied the top positions in the tree-based SHAP rankings, with polycentric distance variables following closely. Distance to Consolacion was the strongest locational signal in both models, and distance to Cebu Business Park also appeared in the top ten for both, confirming that the market responds to multiple active nodes rather than one. Second, property type variables occupied a large share of the top explanatory positions, which means that the same location does not generate the same premium across condominiums, detached houses, and vacant lots. Third, the MCRAI results were mixed. Education, grocery, recreation, and transport fed into the final positive composite, while security, tourism, and retail density carried negative OLS signs and were treated as diagnostic rather than additive categories.

Taken together, these findings suggest that Metro Cebu's residential price surface is not organized around one center with a smooth decline outward. It behaves more like a corridor with several active peaks and overlapping catchments. That reading is consistent with the broader polycentric literature and with the local planning view of Metro Cebu as a system of linked urban clusters rather than a single dominant core \parencite{giuliano1991subcenters,mcmillen2003employment,heikkila1989cbd,jica2015metrocebu}.

\subsection{Polycentric Structure in the SHAP Top Features}\label{polycentric-structure-in-the-shap-top-features}

The SHAP rankings are the clearest evidence that the model learned a polycentric rather than monocentric price structure. In both tree-based models, property-type dummies and bedrooms count occupied the top positions, but the locational variables that followed were consistently polycentric. Distance to Consolacion ranked fourth overall in both the Random Forest and XGBoost SHAP tables. Distance to Cebu Business Park appeared in the top ten for both models. Longitude also ranked among the top features in XGBoost, which fits Metro Cebu's north-south coastal form and the east-west separation created by shoreline, bridges, airport land, and mountain terrain.

This does not mean that the traditional core disappeared. In the OLS baseline, distance to Cebu Business Park still carried a significant negative sign. A better reading is that the old core still matters, but it no longer explains the market by itself. The result fits the polycentric logic of \textcite{giuliano1991subcenters}, the uneven gradient argument of \textcite{mcmillen2003employment}, and the question raised by \textcite{heikkila1989cbd} about whether subcenters can weaken a single-CBD explanation. In Metro Cebu, the market appears layered rather than centerless.

That interpretation also agrees with the JICA roadmap, which treats Metro Cebu as a linked urban system with specialized clusters and corridor growth nodes such as SRP and Naga \parencite{jica2015metrocebu}. Distance to Naga City also appeared in the XGBoost top-ten features even though Naga was outside the training LGU set. Its presence as an external anchor in one of the two models suggests that the southern corridor continues to respond to infrastructure and employment signals that extend beyond administrative training boundaries, which is consistent with the roadmap's treatment of Naga as a corridor anchor.

\subsection{Property Type and the Mactan Island Premium}\label{property-type-and-the-mactan-island-premium}

Property segmentation is part of the main structure of value formation, not a side result. The top SHAP ranks were crowded by property type indicators such as single detached, condominium, and house and lot. This means the model was learning which residential products absorb location advantages more strongly.

This matters for the Mactan side of the market. The results do not support an exaggerated claim that the island dummy alone drives prices. In the retained SHAP summaries, the island indicator itself did not enter the top 10. At the same time, the broader Mactan signal is clear through the strong role of condominium type, the positive OLS coefficient for Lapu-Lapu City, and the repeated importance of distance to the Mactan CBD. The premium is therefore better read as a submarket effect than as a raw island effect. What the market appears to price is the bundle of airport access, bridge-linked connectivity, tourism-adjacent development, and condominium-oriented demand that characterizes the Mactan corridor.

That reading is consistent with Agosto's finding that mobility and livability matter for Cebu land values \parencite{agosto2017}. It also fits the JICA roadmap's view of the tri-city core as functionally linked but internally differentiated \parencite{jica2015metrocebu}. A useful regional comparison is the Bangkok condominium study, which found that transport-related and domestic accessibility factors can explain a large share of condominium value formation in a congested Southeast Asian market \parencite{tochaiwat2023thai}.

\subsection{MCRAI Findings and the Spatial Sorting Interpretation}\label{mcrai-findings-and-the-spatial-sorting-interpretation}

The MCRAI results need a more careful reading than a simple positive-versus-negative split. The final composite retained education, grocery, recreation, and transport because those categories carried positive Stage 1 OLS coefficients in the open-market sample. That result is sensible because these variables describe everyday residential convenience and lower routine household friction.

The negative signs on security, tourism, and retail density should not be read as proof that those categories are inherently harmful. A better interpretation is spatial sorting. In a hedonic framework, households sort across locations according to the bundles of attributes they prefer and can afford \parencite{tiebout1956,bayer2012tiebout}. A variable can therefore carry a negative coefficient not because the amenity itself is undesirable, but because it is concentrated in places with congestion, visitor pressure, commercial intensity, or other traits that push some residential buyers away. The broader capitalization literature supports this conditional reading \parencite{dronyk2017goldilocks,brasington2024fire,yang2016externality,song2004mixed,chenjim2010}.

This is why the Stage 2 MCRAI weight derivation was important. The thesis did not force all categories into one additive score after the Stage 1 signs turned out to be mixed. Instead, it retained the positive household-access categories for the Stage 2 composite and treated the negative categories as diagnostic evidence of sorting and threshold effects. That also lines up with the Metro Manila hedonic report of Tiglao and Rivera and with the Bangkok condominium study, both of which show that accessibility matters, but that the market rewards specific forms of accessibility rather than all forms equally \parencite{tiglaorivera2006,tochaiwat2023thai}.

\subsection{Comparison to Prior Cebu and Manila Hedonic Literature}\label{comparison-to-prior-cebu-and-manila-hedonic-literature}

The findings are broadly consistent with the earlier Cebu literature, but they refine it. Agosto's conclusion that mobility and livability matter is still visible in the present model \parencite{agosto2017}. Transport survived into the final MCRAI composite, and the positive categories mostly describe the daily convenience side of livability. What the current model adds is a stronger spatial reading: this capitalization does not operate evenly across the metropolis, but through several active nodes and property-specific submarkets.

The Metro Manila comparison points in the same direction. Tiglao and Rivera treated accessibility to work and urban opportunities as central to both land valuation and location choice \parencite{tiglaorivera2006}. The Metro Cebu results support that idea, but they also suggest that accessibility in a polycentric region is not a single measure. It is a set of overlapping access logics: access to the old core, access to corridor nodes, access to airport-linked condominium markets, and access to daily neighborhood services. In that sense, Metro Cebu sits between a simple monocentric model and the stronger polycentric claim associated with \textcite{heikkila1989cbd}.

A regional comparison also helps explain why the MCRAI categories did not all move in the same direction. The Bangkok condominium article reported that transport accessibility and related factors could account for 51.69\% of total model weight in that market \parencite{tochaiwat2023thai}. The present study found the same broad importance of accessibility, but in a more selective way because the Metro Cebu model covered several residential property types and a wider mix of corridor conditions.

\subsection{Model Trade-offs}\label{model-trade-offs}

The retained baseline Random Forest model was the best predictive choice under the study's overall evaluation rule, but the findings show why model selection in this thesis could not be treated as a pure one-metric contest. OLS remained useful because it exposed coefficient direction, supported the Stage 1 MCRAI screening, and showed that Cebu Business Park still kept a meaningful role in the price structure. The tree-based models were needed because the market contains non-linear interactions across distance, corridor position, and property type that a single global equation cannot represent well.

The tuned models in Chapter 7 also matter for interpretation. Their weaker out-of-sample performance suggests that the Metro Cebu sample did not reward added complexity by default. The simpler baseline tree models generalized better. That is consistent with recent housing-prediction work that treats tree ensembles as strong tabular learners while still warning that data quality and market heterogeneity remain the real constraints \parencite{nyanda2024,sousa2024}.

\subsection{Limitations}\label{limitations-ch8}

Several limits should stay visible when reading these findings. The model was trained on open-market listings rather than transaction records, so the target variable reflects asking-price behavior as well as market value expectations. The sample was also limited to six LGUs in the training set even though the price surface clearly responds to nearby external nodes such as Naga City. In addition, the MCRAI categories summarize accessibility through carefully designed but still simplified proxies. They cannot fully represent route quality, travel time variability, perceived safety, or institutional quality.

There is also a methodological limit in the way spatial heterogeneity was handled. The tree-based models captured non-linearity and interaction effects, but the thesis did not estimate spatially varying coefficients directly. That matters in a polycentric region because the value of distance to a node is unlikely to be constant across all locations. The polycentric literature suggests exactly this kind of uneven gradient structure \parencite{mcmillen2003employment,heikkila1989cbd}.

Even with these limits, the main interpretation is stable. Metro Cebu's residential valuation pattern is best read as a polycentric market shaped by multiple active centers, property-type segmentation, and selective capitalization of accessibility. Chapter 9 builds on this interpretation by drawing the main conclusions of the thesis, while Chapter 10 turns those conclusions into the study's formal recommendations.
